\subsection{Construire l'intuition avec quatre entités}
Imaginons quatre entités capables d'extraire une ressource de leur
environnement. Chacune dispose d'un stock mobilisable $K_i$ : plus elle
est grande, plus elle extrait, mais chaque unité supplémentaire de stock
rapporte moins que la précédente. La fonction $f_i(K)=A_iK^{\gamma_i}$
décrit cette \emph{technologie d'extraction}. L'apport rejoint le stock,
dont une fraction se dégrade à chaque pas. De nouvelles entités apparaissent
avec la même dotation $K^\circ$ et, dans le régime de référence,
le même exposant de production $\gamma^\circ=1/2$.

Une entité bien dotée peut prêter du stock à une autre, où ce stock produit
davantage à la marge. Ce transfert réel crée simultanément une créance
chez la prêteuse et une dette chez l'emprunteuse. Le \emph{principal} $q$
est le montant nominal inscrit ; le \emph{service des intérêts} $rq$ est
le paiement exigé à chaque pas. Les intérêts déplacent la ressource entre
entités, sans en créer. Le principal n'est pas remboursé progressivement
dans ce modèle : c'est une hypothèse institutionnelle forte.

\begin{figure}[htbp]\centering
\begin{tikzpicture}[>=Stealth,font=\small,
 ent/.style={circle,draw=blue!60!black,fill=blue!7,minimum size=15mm,align=center},
 lab/.style={font=\scriptsize,fill=white,inner sep=2pt}]
\draw[rounded corners,fill=gray!3,draw=gray!50] (-1.5,-1.7) rectangle (9.5,2.3);
\node[ent] (a) at (0,0.7) {$e_1$\\$K_1=16$};
\node[ent] (b) at (4,0.7) {$e_2$\\$K_2=4$};
\node[ent] (c) at (7.5,1.2) {$e_3$\\$K_3=9$};
\node[ent] (d) at (7.5,-0.8) {$e_4$\\$K_4=1$};
\draw[->,very thick,blue!70!black] (a) to[bend left=18]
 node[lab,above] {prêt initial $q=6$} (b);
\draw[->,dashed,thick,orange!80!black] (b) to[bend left=22]
 node[lab,below] {intérêts $rq$, à chaque pas} (a);
\draw[->,green!45!black,thick] (-.9,2.8)--(a);
\node[lab] at (2.4,2.8) {apport extrait de l'environnement : $f_i(K_i)$};
\draw[->,green!45!black,thick] (7.5,2.8)--(c);
\draw[->,red!70!black] (d)--(9.8,-1.4) node[right,font=\scriptsize] {érosion};
\node[align=center,font=\scriptsize] at (2.7,-1.25)
 {Après le prêt : $K_1=K_2=10$ ; $C_1=6$, $D_2=6$.\\
 Les valeurs nettes restent $NW_1=16$ et $NW_2=4$.};
\end{tikzpicture}
\caption{\textbf{Stock réel et engagement comptable ne sont pas la même chose.}
Exemple fictif, inspiré du schéma pédagogique de la présentation du stage.
Les valeurs initiales sont choisies pour expliquer le mécanisme, non extraites
d'une simulation. Les entités 3 et 4 montrent que toutes n'ont pas à participer
au même échange. Les apports et l'érosion ne sont dessinés que sur quelques
entités pour alléger le schéma ; les règles s'appliquent à toutes.
Le réservoir extérieur n'est pas simulé.}
\label{fig:exemple}
\end{figure}

Le prêt conserve ainsi la valeur nette instantanée de chaque partenaire,
mais change sa capacité productive et ses engagements futurs. Une entité
peut ensuite ne plus pouvoir honorer son service, ou devoir plus qu'elle
ne possède. Sa disparition efface les créances détenues sur elle : ses
prêteuses peuvent à leur tour devenir insolvables. Le couplage entre stocks
réels, engagements persistants et pertes propagées est le mécanisme étudié.
